\chapter{Decision Extraction}
\label{ch:decision-extraction}

\textit{The extraction pipeline is the heart of \sysname{}.  It takes a
free-form conversation---possibly thousands of tokens of rambling
developer--AI dialogue---and distills it into a structured decision trace
with six fields.  This chapter walks through every stage of that pipeline,
from prompt design to JSON repair to caching.}


% ─────────────────────────────────────────────────────────────────────
\section{The Pipeline}
\label{sec:extraction-pipeline}

The extraction service lives in
\codefile{apps/api/services/extractor.py}.  At a high level the pipeline
is five stages:

\begin{enumerate}
  \item \textbf{Input.}  A block of conversation text arrives---either from a
        completed capture session (Chapter~\ref{ch:your-first-session}) or
        from a bulk JSON import.
  \item \textbf{Prompt construction.}  The text is injected into a
        few-shot extraction prompt that tells the LLM what a ``decision''
        is and what fields to return.
  \item \textbf{LLM call.}  The prompt is sent to the configured LLM
        (default: NVIDIA Nemotron 49B) with \ic{temperature=0.3} and
        \ic{max\_tokens=4096}.
  \item \textbf{JSON parse \& repair.}  The raw response is cleaned
        (thinking tags stripped, truncated JSON repaired) and parsed into a
        Python list of dicts.
  \item \textbf{Validation \& defaults.}  Missing fields are filled with
        safe defaults, and each dict is validated against the
        \ic{DecisionCreate} schema.
\end{enumerate}

\begin{keyinsight}
The pipeline is intentionally \emph{stateless}.  It takes text in and
returns structured JSON out.  Persistence to Neo4j, PostgreSQL, and Redis
happens in a separate \ic{save\_decision()} step---which means you can
test extraction in isolation without touching any database.
\end{keyinsight}


% ─────────────────────────────────────────────────────────────────────
\section{The Extraction Prompt}
\label{sec:extraction-prompt}

The core prompt (\ic{DECISION\_EXTRACTION\_PROMPT}) is a few-shot,
chain-of-thought template.  Its structure has three parts:

\paragraph{1.\ Definition of ``decision.''}
The prompt defines four categories that the LLM should recognize:

\begin{itemize}
  \item \textbf{Explicit decisions:} ``Should we use X or Y?  Let's use X
        because\ldots''
  \item \textbf{Implicit decisions:} ``Let's use X for this'' (no
        stated alternatives).
  \item \textbf{Technical choices:} Framework selections, architecture
        patterns, tool adoptions.
  \item \textbf{Implementation strategies:} How to solve a problem,
        approach to take.
\end{itemize}

\paragraph{2.\ Required fields.}
Each decision must include six fields:

\begin{table}[ht]
\centering
\caption{The six fields of a decision trace.}
\label{tab:trace-fields}
\begin{tabularx}{\textwidth}{l X}
\toprule
\textbf{Field} & \textbf{Description} \\
\midrule
\ic{trigger}    & What prompted the decision (problem, requirement, or question). \\
\ic{context}    & Background information---constraints, requirements, team situation. \\
\ic{options}    & Alternatives considered (may be a single-element list). \\
\ic{decision}   & A clear statement of what was chosen. \\
\ic{rationale}  & Why this choice was made, or ``Not explicitly stated.'' \\
\ic{confidence} & 0.0--1.0 score reflecting how clear the decision was. \\
\bottomrule
\end{tabularx}
\end{table}

\paragraph{3.\ Output format.}
The prompt closes with: ``Return ONLY valid JSON, no markdown code blocks
or explanation.''  The expected output is a JSON array of decision objects,
or an empty array \ic{[]} if no decisions are found.

The prompt also includes four few-shot examples (single decision, multiple
decisions, implicit decision, and no-decision conversation) to anchor the
LLM's behavior.  An abbreviated version:

\begin{lstlisting}[style=python, caption={Abbreviated extraction prompt (see extractor.py for the full version).}]
DECISION_EXTRACTION_PROMPT = """Analyze this conversation
and extract any technical decisions made.

## What constitutes a decision?
A decision is a choice that affects the project direction,
architecture, or implementation. This includes:
- Explicit decisions: "Should we use X or Y? ..."
- Implicit decisions: "Let's use X for this"
- Technical choices: Framework, architecture, tools
- Implementation strategies: How to solve a problem

## Examples
### Example 1: Single clear decision
Conversation: "We need to pick a database. ..."
Output: [{"trigger": "...", "decision": "...", ...}]

### Example 4: No decisions (just discussion)
Output: []

## Conversation to analyze:
{conversation_text}

Return ONLY valid JSON, no markdown code blocks."""
\end{lstlisting}

\sysname{} also ships specialized prompts for architecture, technology,
and process decisions (\ic{ARCHITECTURE\_DECISION\_PROMPT},
\ic{TECHNOLOGY\_DECISION\_PROMPT}, \ic{PROCESS\_DECISION\_PROMPT}).
These narrow the LLM's focus when the decision type is known in advance.


% ─────────────────────────────────────────────────────────────────────
\section{JSON Repair}
\label{sec:json-repair}

LLM responses are not always well-formed JSON.  Two problems are common:

\paragraph{Thinking tags.}
NVIDIA Nemotron generates \ic{<think>...</think>} blocks as part of its
chain-of-thought reasoning.  These consume output tokens and are not valid
JSON.  The system strips them before parsing.

\paragraph{Truncated output.}
When the combined thinking + JSON exceeds \ic{max\_tokens}, the response
is cut off mid-object.  The JSON repair strategy in
\codefile{utils/json\_extraction.py} tries four strategies in order:

\begin{enumerate}
  \item Parse as pure JSON.
  \item Extract from \ic{\`{}\`{}\`{}json} code blocks.
  \item Extract from untyped \ic{\`{}\`{}\`{}} code blocks.
  \item Regex fallback: find the last complete \ic{\}} followed by
        \ic{]} and parse the substring.
\end{enumerate}

\begin{warning}
Always use \ic{max\_tokens=4096}, not the default 2000.  Thinking tags
routinely consume 1000--2000 tokens before the model starts producing
JSON.  With \ic{max\_tokens=2000}, roughly 14.5\% of extractions fail
due to truncation.
\end{warning}


% ─────────────────────────────────────────────────────────────────────
\section{Prompt Versioning}
\label{sec:prompt-versioning}

LLM extraction responses are cached in Redis to avoid redundant API calls.
The cache key includes the prompt version:

\begin{lstlisting}[style=python, numbers=none]
# Format: llm:{version}:{type}:{md5(text)}
key = f"llm:{version}:extract:{text_hash}"
\end{lstlisting}

The version string is controlled by
\confkey{LLM\_EXTRACTION\_PROMPT\_VERSION} in
\codefile{apps/api/config.py} (default: \ic{"v1"}).  When you modify the
extraction prompt, bump this value.  All cached responses whose keys
contain the old version will naturally miss, forcing re-extraction with the
new prompt.

\begin{lstlisting}[style=python, numbers=none, caption={The cache key generation method.}]
def _get_cache_key(self, text: str, extraction_type: str) -> str:
    text_hash = hashlib.md5(text.encode("utf-8")).hexdigest()
    version = self._settings.llm_extraction_prompt_version
    return f"llm:{version}:{extraction_type}:{text_hash}"
\end{lstlisting}

Cached responses have a 24-hour TTL (\confkey{llm\_cache\_ttl} =
86\,400\,s).  Entity resolution results are cached separately with a
5-minute TTL.


% ─────────────────────────────────────────────────────────────────────
\section{Default Values}
\label{sec:default-values}

Not every LLM response includes all six fields.  The
\ic{apply\_decision\_defaults()} function fills in safe values for anything
missing or empty:

\begin{lstlisting}[style=python, caption={Default values for incomplete extraction results.}]
DEFAULT_DECISION_FIELDS = {
    "confidence": 0.5,
    "context":    "",
    "rationale":  "",
    "options":    [],
    "trigger":    "Unknown trigger",
    "decision":   "",
}
\end{lstlisting}

The function also preserves extra fields not in the defaults dict, so
specialized prompts that return \ic{decision\_type} or other custom fields
are not silently dropped.

\begin{tryit}
Run a single extraction manually to see defaults in action:
\begin{lstlisting}[style=bash, numbers=none]
curl -X POST http://localhost:8000/api/ingest \
  -H "Content-Type: application/json" \
  -d '{"text": "Let us use Redis for caching."}'
\end{lstlisting}
The response will show \ic{confidence: 0.85} (from the LLM) but
\ic{context: ""} and \ic{options: ["Redis"]} (both filled by the model),
demonstrating that the LLM usually provides most fields when the
conversation is clear.
\end{tryit}


% ─────────────────────────────────────────────────────────────────────
\section{Gotchas}
\label{sec:extraction-gotchas}

\begin{warning}
\textbf{Prompt sanitizer blocks synthetic conversations.}
The LLM client includes a prompt injection detector that flags
``USER:'' and ``ASSISTANT:'' labels as potential attacks.  Synthetic
conversations (used in evaluation) naturally contain these markers.
Always pass \ic{sanitize\_input=False} when calling the LLM on
synthetic or imported conversation text.
\end{warning}

\begin{warning}
\textbf{Multiple extraction runs accumulate.}
The \ic{save\_decision()} method \emph{creates} new Neo4j nodes---it
does not upsert.  If you run the extraction pipeline twice on the same
conversations, you get duplicate decision traces.  Always wipe Neo4j
(\ic{MATCH (n) DETACH DELETE n}) before starting a new evaluation run.
\end{warning}

\begin{warning}
\textbf{True extraction yield is approximately 1.9 decisions per
conversation} with an 85.5\% success rate across 200 synthetic
conversations.  If you see significantly higher numbers, you are likely
looking at accumulated results from multiple runs.
\end{warning}
